
\documentclass[a4paper, 12]{article}


\usepackage{amsmath,amssymb} 
\usepackage[french]{babel}
                                                                                                                               
\newtheorem{ex}{Exemple}

\newcommand\set[1]{\ensuremath{\{#1\}}}



\begin{document}

\title{Partitionner les sommes des valeurs des cartes \\avec des quantiles}
\date{}
\maketitle


\section{Partition des sommes des valeurs des cartes dans Kimind}

Les valeurs intéressantes sont les nombres de façons de distribuer les sommes sur 4 cartes. Pour une somme donnée, plus ce nombre est grand plus le problème est difficile. Pour qualifier une somme nous avons besoin de deux séries statistiques.

\subsection{Série sur le nombre de façons de distribuer les valeurs des cartes pour des sommes données}

On calcule le nombre de façons de distribuer les valeurs sur 4 cartes pour chaque somme  :

\begin{itemize}
	\item $1$ distribution pour la somme $4$, 
	\item $1$ distribution pour la somme $40$,
	\item $4$ distributions pour la somme $5$,
	\item $4$ distributions pour la somme $39$,
	\item $\cdots$,
	\item $660$ distributions pour la somme $21$,
	\item $660$ distributions  pour la somme $23$,
	\item $670$ distributions pour la somme  $22$.
\end{itemize}
(voir annexe \ref{distribution1}.)

On a donc la série d'effectifs suivante où les valeurs sont les nombres de distributions  :
\set{e(1,2), e(4,2), \cdots, e(660, 2), e(670,1)} qui correspondent à la série 
$D = \set{1, 1, 4, 4, \cdots,660,660,670}$ ayant $37$ valeurs. Noter que les sommes n'apparaissent pas dans cette série.

On calcule la fonction de répartition $F(d_i)$ pour chaque valeur \underline{distincte} $d_i$ de $D$ :
\begin{itemize}
\item $F(1) = 2/37$ : proportion des distributions dans $D$ dont la valeur est inférieure ou égale à $1$,
\item $F(4) = 4/37$ ($4$ distributions de $D$ sont inférieures ou égales à $4$),
\item $\cdots$,
\item $F(670) = 37/37$ ($37$ distributions de $D$ sont inférieures ou égales à $670$).
\end{itemize}

Après on calcule les quintiles :
\begin{itemize}
\item $Q_1 = Q(0.2)$ : on cherche le plus petit $d_i$ de $D$ tel que $F(x_i) \ge 0.2$,
\item $Q_2 = Q(0.4)$ : on cherche le plus petit $d_i$ de $d$ tel que $F(x_i) \ge 0.4$,
\item $Q_3 = Q(0.6)$ : on cherche le plus petit $d_i$ de $D$ tel que $F(x_i) \ge 0.6$,
\item $Q_4 = Q(0.8)$ : on cherche le plus petit $d_i$ de $D$ tel que $F(x_i) \ge 0.8$.
\end{itemize}

puis on créé le parties :
\begin{itemize}
	\item $D_1$ = $\set{d_i \in D   | \ d_i \le Q_1}$,
	\item $D_2$ = $\set{d_i \in D   | \ Q_1 < d_i \le Q_2}$,
	\item $D_3$ = $\set{d_i \in D   | \ Q_2 < d_i \le Q_3}$,
	\item $D_4$ = $\set{d_i \in D   | \ Q_3 < d_i \le Q_4}$,
	\item $D_5$ = $\set{d_i \in D   |  d_i >  Q_4}$,
\end{itemize}.

On attribue aux distributions de $D_1$ le qualificatif \emph{tres$\_$facile}, pour $D_2$ \emph{facile}, pour $D_3$ \emph{moyen}, pour $D_4$ \emph{difficile} pour $D_5$ \emph{tres$\_$difficile}. 

\subsection{Série sur les sommes}
Nous calculons les \emph{effectifs des sommes} $ S = \set{e(4,1), e(5,4), \cdots, e(39,4), e(40,1)}$ contenant les couples $(s_i, d_i)$ où $s_i$ est une somme et $d_i$ le nombre de façons de distribuer cette somme sur 4 cartes (voir annexe \ref{somme1}). 

\subsection{Partition de la série sur les sommes}

L'idée est d'utiliser la partition de $D$ pour partitionner $S$ en 4 ensembles $S_1$ (\emph{sommes$\_$tres$\_$faciles}), $S_2$ (\emph{sommes$\_$faciles}), $S_3$ (\emph{sommes$\_$moyennes}),\\  $S_4$ (\emph{sommes$\_$difficiles}) et 
$S_5$ (\emph{sommes$\_$tres$\_$difficiles}).

Pour chaque somme-effectif $e(s_i, d_i)$ de $S$ on regarde dans quelle partie $D_j$ de $D$ se trouve $d_i$ et l'on place $s_i$ dans la partie $S_j$ de $S$. 

Par exemple, on souhaite connaître la difficulté associée à la somme-effectif $e(39, 4)$ de $S$. On constate que $4$ (nombre de façons de distribuer $39$) est dans la partie $D_1$ de $D$, on met donc la somme $39$ dans $S_1$ comprenant les sommes \emph{tres$\_faciles$}.


\subsection{Prédicats disponibles}

Le prédicat suivant  de \emph{fonctions$\_$somme.pl} retourne les deux séries nécessaires :
\begin{verbatim}
effectifs_sommes_distributions_positions(Somme_min, Somme_max, Nb_cartes, 
                                         Min, Max, S, D).
\end{verbatim}
qui constitue 2 séries statistiques (avec les  effectifs) :
\begin{itemize}
\item la série \emph{S} contient les sommes  (voir annexe \ref{somme1}).
	  
\item la série \emph{D}  contient les nombres de distributions des sommes (voir annexe \ref{distribution1}).
\end{itemize}
Par exemple l'appel 
\begin{verbatim}
effectifs_sommes_distributions_positions(4, 40, 4, 1, 10, S, D)
\end{verbatim}
retourne les effectifs des sommes $S$ et les effectifs des nombres de distributions des sommes $D$, pour des sommes entre $4$ et $40$, avec $4$ cartes, chaque carte ayant une valeur minimale de $1$ 
et maximale de $10$.

Par ailleurs, si l’on ne veut pas tenir compte des positions on appelle le prédicat 
\begin{verbatim}
effectifs_distributions_sommes(Somme_min, Somme_max, Nb_cartes, 
                               Min, Max, S, D).
\end{verbatim}


Le prédicat \emph{difficultes$\_$sommes$\_$positions(Nb$\_$difficultes, Partition$\_$sommes)} retourne la partition des sommes en fonction de leurs difficultes). Les sommes de $4$ à $40$ sont considérées et les difficultés possibles sont \emph{$[$tres$\_$facile, facile, moyen, difficile, tres$\_$difficile$]$}  (annexe \ref{annexe45}) soit \emph{$[$facile, moyen, difficile$]$} (annexe \ref{annexe46}).

Le prédicat \emph{difficultes$\_$sommes(Nb$\_$difficultes, Partition$\_$sommes)}  partitionne les sommes sans tenir compte des positions
(voir annexes \ref{annexe47} et \ref{annexe48}).

\section{Annexes}

\subsection{Effectifs des sommes (en tenant compte des positions)}
\label{somme1}

il y a 10000 sommes dans la série, voici les effectifs (somme, nbOccurrences) :

\begin{verbatim}
e(4,1)	% la somme 4 a 1 distribution
e(5,4)      % la somme 5 a 4 distributions
e(6,10)
e(7,20)
e(8,35)
e(9,56)
e(10,84)
e(11,120)
e(12,165)
e(13,220)
e(14,282)
e(15,348)
e(16,415)
e(17,480)
e(18,540)
e(19,592)
e(20,633)
e(21,660)
e(22,670)
e(23,660)
e(24,633)
e(25,592)
e(26,540)
e(27,480)
e(28,415)
e(29,348)
e(30,282)
e(31,220)
e(32,165)
e(33,120)
e(34,84)
e(35,56)
e(36,35)
e(37,20)
e(38,10)
e(39,4)
e(40,1)   % la somme 40 a 1 distribution
\end{verbatim}


\subsection{Effectifs des nombres de distributions des sommes en tenant compte des positions}
\label{distribution1}
Il y a 37 nombres de distributions dans la série voici les effectifs (nbDistribution, NbOccurrences)

\begin{verbatim}
   e(1,2),	  %  2 sommes distinctes ont 1 distribution possible
   e(4,2),    % 2 sommes ont 4 distributions (la valeur 4 apparait 2 fois dans la série)
   e(10,2),  % 10 distributions possibles pour 2 sommes
   e(20,2),
   e(35,2),
   e(56,2),
   e(84,2),
   e(120,2),
   e(165,2),
   e(220,2),
   e(282,2),
   e(348,2),
   e(415,2),
   e(480,2),
   e(540,2),
   e(592,2),
   e(633,2),
   e(660,2), 
   e(670,1)   % 1 somme a 670 distributions
\end{verbatim}



\subsection{Effectif des sommes (sans tenir compte des positions)}

il y a 715 sommes dans la série voici les effectifs (somme, nbOccurrences) :

\begin{verbatim}
   e(4,1),		%  la somme 4 a 1 distribution (elle n'apparait qu'une fois dans la série)
   e(5,1)		% la somme 5 a 1 distribution
   e(6,2),		% la somme 6 a 2 distributions
   e(7,3),
   e(8,5),
   e(9,6),
   e(10,9),
   e(11,11),
   e(12,15),
   e(13,18),
   e(14,22),
   e(15,25),
   e(16,30),
   e(17,32),
   e(18,36),
   e(19,38),
   e(20,41),
   e(21,41),
   e(22,43),
   e(23,41),
   e(24,41),
   e(25,38),
   e(26,36),
   e(27,32),
   e(28,30),
   e(29,25),
   e(30,22),
   e(31,18),
   e(32,15),
   e(33,11),
   e(34,9),
   e(35,6),
   e(36,5),
   e(37,3),
   e(38,2),
   e(39,1),	% la somme 39 a 1 distribution
   e(40,1)         % la somme 40 a 1 distribution
\end{verbatim}

\subsection{Effectif des nombres de distributions des sommes (sans tenir compte des positions)}

il y a $37$ éléments dans la série, voici les effectifs (nbDistribution, NbOccurrences) :

\begin{verbatim}
   e(1,4),		% 4 sommes ont 1 distributions
   e(2,2),
   e(3,2),
   e(5,2),
   e(6,2),
   e(9,2),
   e(11,2),
   e(15,2),
   e(18,2),
   e(22,2),
   e(25,2),
   e(30,2),
   e(32,2),
   e(36,2),
   e(38,2),
   e(41,4),
   e(43,1)     % 1 somme a 43 distributions
\end{verbatim}


\subsection{Difficultés  avec positions des sommes avec 5 classes de difficultés}
\label{annexe45}
\begin{verbatim}
q(very_easy,[e(4,1),e(5,4),e(6,10),e(7,20),e(37,20),e(38,10),e(39,4),e(40,1)])

q(easy,[e(8,35),e(9,56),e(10,84),e(11,120),e(33,120),e(34,84),e(35,56),e(36,35)])

q(average,[e(12,165),e(13,220),e(14,282),e(15,348),e(29,348),e(30,282),
           e(31,220),e(32,165)])

q(difficult,[e(16,415),e(17,480),e(18,540),e(26,540),e(27,480),e(28,415)])

q(very_difficult,[e(19,592),e(20,633),e(21,660),e(22,670),e(23,660),
                  e(24,633),e(25,592)])
\end{verbatim}


\subsection{Difficultés  avec positions des sommes avec  3 classes de difficultés}
\label{annexe46}
\begin{verbatim}
q(easy,[e(4,1),e(5,4),e(6,10),e(7,20),e(8,35),e(9,56),e(10,84),e(34,84),
	       e(35,56),e(36,35), e(37,20),e(38,10),e(39,4),e(40,1)])
   
q(average,[e(11,120),e(12,165),e(13,220),e(14,282),e(15,348),e(16,415),
           e(28,415),e(29,348),e(30,282),e(31,220),e(32,165),e(33,120)])
   
q(difficult,[e(17,480),e(18,540),e(19,592),e(20,633),e(21,660),
             e(22,670),e(23,660), (24,633),e(25,592),e(26,540),e(27,480)])
\end{verbatim}


\subsection{Difficultés  sans positions des sommes avec  5 classes de difficultés}
\label{annexe47}
\begin{verbatim}
q(very_easy,[e(4,1),e(5,1),e(6,2),e(7,3),e(37,3),e(38,2),e(39,1),e(40,1)])

q(easy,[e(8,5),e(9,6),e(10,9),e(11,11),e(33,11),e(34,9),e(35,6),e(36,5)])

q(average,[e(12,15),e(13,18),e(14,22),e(15,25),e(29,25),e(30,22),e(31,18),e(32,15)])

q(difficult,[e(16,30),e(17,32),e(18,36),e(26,36),e(27,32),e(28,30)])

q(very_difficult,[e(19,38),e(20,41),e(21,41),e(22,43),e(23,41),e(24,41),e(25,38)])
\end{verbatim}

\subsection{Difficultés  sans positions des sommes avec  3 classes de difficultés}
\label{annexe48}
\begin{verbatim}
q(easy,[e(4,1),e(5,1),e(6,2),e(7,3),e(8,5),e(9,6),e(10,9),e(34,9),e(35,6),
        e(36,5),e(37,3),e(38,2),e(39,1),e(40,1)])

q(average,[e(11,11),e(12,15),e(13,18),e(14,22),e(15,25),e(16,30),e(28,30),
           e(29,25),e(30,22),e(31,18),e(32,15),e(33,11)])

q(difficult,[e(17,32),e(18,36),e(19,38),e(20,41),e(21,41),e(22,43),e(23,41),
             e(24,41),e(25,38),e(26,36),e(27,32)])
\end{verbatim}

\end{document}